\documentclass[10pt,a4paper]{amsart}
\usepackage[margin=19mm]{geometry}
\usepackage{amsmath,amssymb,mathtools}
\usepackage[scr=rsfs,cal=euler]{mathalfa}
\usepackage{xfrac}
\usepackage[unicode,hidelinks,bookmarks=false]{hyperref}
\newcommand{\ie}{i.e.\ }
\newcommand{\cf}{cf.\ }
\newcommand{\resp}{respectively\ }
\newcommand{\Subsection}{\subsection}
\providecommand{\nobreakdash}{\nobreak\hbox{-}}
\setlength{\emergencystretch}{2em}
\multlinegap0pt
\usepackage{enumerate}
\usepackage{float}
\usepackage{amssymb}
\usepackage{tikz}
\newtheorem{lemma}{Lemma}
\title{Settling the total domination-annihilation conjecture for graphs with minimum degree two}
\author{Marko Jakovac$^{a,b,}$}
\date{Source: arXiv:2609.10795v2 (CC BY 4.0)}
\begin{document}
\maketitle

\noindent\emph{Source and licence.} Settling the total domination-annihilation conjecture for graphs with minimum degree two. Version arXiv:2609.10795v2 (CC BY 4.0), distributed by arXiv under the Creative Commons Attribution 4.0 International licence, http://creativecommons.org/licenses/by/4.0/, as recorded in the arXiv licence metadata for this exact version and shown on the abstract page https://arxiv.org/abs/2609.10795v2. Full licence notice: \texttt{licenses/arxiv-2609.10795-CC-BY-4.0.txt}.\par\smallskip

\noindent\textit{Editorial scope.} Counted unit: the article's lemma lem:d*3 (source0.tex lines 153--156). The article's complete original proof of it is delivered with it (source0.tex lines 158--204, one \texttt{proof} environment of 47 lines). No mathematics is written, repaired or re-derived: the statement and the proof are the article's own lines, in the article's own notation, and no retained line is substituted or re-flowed.  Retained uncounted as the source-local context the proof needs: lines 135--137.  Nothing was omitted from any retained span, so the package carries no omission record.  No retained line contains a citation, so the wrapper carries no bibliography and no bibliography entry.  No retained line contains a figure environment, an image-inclusion command or drawing code, so no image is delivered and the package declares no asset dependency.  Editorial formatting only, in addition: an AMS \texttt{amsart} preamble that restates the article's own declarations and macros used by the retained lines, namely the environments \texttt{lemma} on separate counters, together with the standard packages those lines need.  The retained environments print in this document as Lemma 1 in order of appearance, whereas the article prints the same items with the numbering of its own full document.  The complete native source of the article as distributed in the arXiv e-print for this exact version is bundled unchanged as \texttt{sources/209951/source0.tex}.
\begin{equation}\label{eq:p-s}
3p(G)\le |S|.
\end{equation}

\begin{lemma}\label{lem:d*3}
Let $G$ be a connected graph of order $n$ with $\delta(G)\ge 2$ and $\Delta(G)\ge 3$. If $d^*(G)\ge3$, then
$$a(G)\ge \left\lfloor\frac{n(G)+p(G)}2\right\rfloor.$$
\end{lemma}

\begin{proof}
Let $S$ be the set of degree-two vertices and $L=V(G)\setminus S$ the set of large vertices. Define $s=|S|$, $r=|L|$, and let
$$\ell=\sum_{v\in L}d(v)$$
be the degree sum of the large vertices. Since every vertex in $S$ has degree $2$, we have
\begin{equation}\label{eq:m-l}
2m(G)=2s+\ell, \qquad m(G)=s+\frac{\ell}{2}.
\end{equation}
Moreover,
\begin{equation}\label{eq:l-3r}
 \ell\ge3r,
\end{equation}
because every vertex of $L$ has degree at least $3$.

The assumption $d^*(G)\ge 3$ implies that all $s$ degree-two vertices occur among the first $a(G)$ vertices in the nondecreasing degree sequence. In particular, they are contained in an optimal annihilation set. Therefore $2s\le m(G)$. Using \eqref{eq:m-l}, we get
\begin{equation}\label{eq:l-2s}
\ell \ge 2s.
\end{equation}

We now order the degrees of the large vertices as
$$e_1\le e_2\le\cdots\le e_r, \qquad \sum_{i=1}^{r}e_i=\ell.$$
Set
$$t=\left\lfloor \frac{r}{2}-\frac{r\cdot s}{\ell}\right\rfloor.$$
By \eqref{eq:l-2s}, we have $t\ge 0$. Suppose first that $t\ge 1$. Since $e_1,\ldots,e_t$ are the $t$ smallest degrees, their average is at most the average of all $r$ degrees. Hence,
$$\frac{1}{t}\sum_{i=1}^{t}e_i\le\frac{\ell}{r},$$
and therefore
$$\sum_{i=1}^{t}e_i\le\frac{t\cdot\ell}{r}.$$
If $t=0$, the sum is considered zero and hence, in both cases,
$$\sum_{i=1}^{t}e_i\le\frac{t\cdot\ell}{r}.$$
Consequently,
$$2s+\sum_{i=1}^{t}e_i \le 2s+\frac{t \cdot \ell}{r} \le 2s+\left(\frac{r}{2}-\frac{r \cdot s}{\ell}\right)\frac{\ell}{r}=s+\frac{\ell}{2}=m(G),$$
where the last equality follows from \eqref{eq:m-l}. Hence, the set consisting of all $s$ vertices of degree two together with the $t$ large vertices of smallest degrees is an annihilation set. Therefore,
\begin{equation}\label{eq:a-st}
a(G)\ge s+t =\left\lfloor s+\frac r2-\frac{r \cdot s}{\ell}\right\rfloor.
\end{equation}

It remains to estimate the expression inside the floor. Using $n(G)=s+r$, we obtain
$$s+\frac r2-\frac{r \cdot s}{\ell}=\frac{s+r}{2}+\frac{s}{2}\left(1-\frac{2r}{\ell}\right)=\frac{n(G)}{2}+\frac{s}{2}\left(1-\frac{2r}{\ell}\right).$$
By \eqref{eq:l-3r},
$$1-\frac{2r}{\ell}\ge 1-\frac{2r}{3r}=\frac{1}{3},$$
and therefore
$$s+\frac r2-\frac{r \cdot s}{\ell}\ge \frac {n(G)}{2}+\frac{s}{6}.$$
Finally, \eqref{eq:p-s} gives $\frac{s}{6}\ge \frac{p(G)}{2}$. Thus
$$s+\frac{r}{2}-\frac{r \cdot s}{\ell} \ge \frac{n(G)+p(G)}{2}.$$
Together with \eqref{eq:a-st}, this yields
$$a(G)\ge\left\lfloor\frac{n+p(G)}2\right\rfloor,$$
as required.
\end{proof}

\end{document}
